%% ---------------------------------------------------------------------
%% 01 - GIỚI THIỆU
%% Đặt vấn đề, khoảng trống nghiên cứu, mục tiêu và đóng góp của bài báo.
%% Khoảng 2-4 đoạn; trích dẫn bằng \cite{...} (biblatex, Vancouver).
%% ---------------------------------------------------------------------

\section{Giới thiệu}

Lập kế hoạch đường đi là một chức năng cốt lõi của robot di động tự hành
trong kho vận, sản xuất, cứu hộ và dịch vụ đô thị. Robot phải tìm được
quỹ đạo không va chạm từ điểm xuất phát đến đích trong thời gian đủ ngắn,
đồng thời hạn chế tối đa việc dừng lại để tính toán lại khi môi trường
thay đổi \cite{karnik2021}.

Các bộ lập kế hoạch dựa trên lấy mẫu như RRT và RRT* có ưu điểm là không
phụ thuộc vào cấu trúc của không gian cấu hình, nhưng chi phí lấy mẫu
tăng nhanh khi số chiều hoặc mật độ vật cản tăng. Ngược lại, các phương
pháp dựa trên biểu diễn hình học như biểu đồ Voronoi cho lời giải có bảo
đảm về khoảng cách an toàn nhưng khó thích nghi với vật cản chuyển động
\cite{bhattacharya2017}.

\emph{(Phần này là văn bản mẫu. Thay bằng nội dung thật: nêu bối cảnh,
khoảng trống nghiên cứu, câu hỏi nghiên cứu và ba đóng góp chính của bài
báo, mỗi đóng góp một câu.)}

\subsection{Đóng góp của bài báo}

Bài báo đề xuất một khung lai ghép ba lớp; các đóng góp chính gồm:
(i)~biểu diễn miền lấy mẫu bằng Voronoi có trọng số kết hợp
$\beta$-complexes; (ii)~cơ chế hướng dẫn lấy mẫu cho RRT* dựa trên cấu
trúc hình học đó; và (iii)~bộ giám sát logic mờ bảo đảm an toàn khi vật
cản di chuyển.
